% sec_06_geometry
\section{Geometry: projection, occlusion, curvature}\label{sec:geom}

With the transmission coefficient fixed at its normal-incidence value $T_0$
(\cref{sec:fresnel}), the absorbed power density of \cref{sec:surface-law}
depends only on the body shape and the incidence direction. The angular factor
\begin{equation}\label{eq:geom-f}
  f(\rr) = \ReLU\!\bigl[\nhat(\rr)\cdot(-\khat)\bigr]
\end{equation}
is purely geometric. It carries the surface-normal field $\nhat(\rr)$ and the
wave direction $\khat$ and nothing else, so millimetre-wave dosimetry reduces to
a problem in integral geometry. This section sets out three pieces of that
geometry: the visibility factor that makes the law non-local on a non-convex
body, the projected-area integral that fixes total absorbed power, and the
principal curvatures that the diffractive gate of \cref{sec:fock-local} needs.

\subsection{Visibility and occlusion}\label{sec:geom-occlusion}

A human body is not convex. Concavities (the armpit, the gap between the legs,
the region behind an ear) let one part of the body shadow another, so the local
value of $\Sab$ is not determined by $\nhat(\rr)$ alone.

\begin{definition}[Occlusion factor]\label{def:occlusion}
For a surface point $\rr$ and an incidence direction $\khat$, the occlusion
factor is
\begin{equation}\label{eq:O-def}
  O(\rr,\khat) = \begin{cases}
    1 & \text{if $\rr$ is visible along $-\khat$ (no other body part blocks it),}\\
    0 & \text{otherwise.}
  \end{cases}
\end{equation}
\end{definition}

The local law becomes
\begin{equation}\label{eq:Sab-occluded}
  \Sab(\rr) = \Sinc\,T_0\,\ReLU\!\bigl[\nhat(\rr)\cdot(-\khat)\bigr]\,O(\rr,\khat).
\end{equation}
Computing $O$ is one ray cast per surface element against the body mesh, the same
operation a renderer performs for hard shadows. For a convex body $O\equiv 1$ and
\cref{eq:Sab-occluded} collapses back to the bare surface law.

The binary factor in \cref{def:occlusion} is the geometric-optics content of
shadowing. It is exact in the limit $\lambda\to 0$ and discontinuous across the
shadow edge. The two smooth, wave-correct refinements of $O$ are the subject of
the next two sections. A body shadowing itself across a local convex edge (its
own terminator) is the Fock gate of \cref{sec:fock-local}. One body part
shadowing another across a finite gap (a hand over the cheek) is the distal
angular-clearance gate of \cref{sec:cmetric}. Both reduce to \cref{eq:O-def} as
$\lambda\to 0$.

\begin{definition}[Exposure fraction]\label{def:eta}
The exposure fraction at $\rr$ is the cosine-weighted fraction of the hemisphere
from which $\rr$ is visible,
\begin{equation}\label{eq:eta-def}
  \eta(\rr) = \frac{1}{\pi}\int_{S^2}
  \ReLU\!\bigl[\nhat(\rr)\cdot(-\khat)\bigr]\,O(\rr,\khat)\,\diff\Omega(\khat).
\end{equation}
\end{definition}

For convex bodies $\eta\equiv 1$. For non-convex bodies $\eta<1$ in concavities
and $\eta=1$ on exposed regions. The exposure fraction is mathematically
identical to the ambient-occlusion factor of computer
graphics~\cite{Zhukov1998}, which imports decades of GPU ray-casting and
bounding-volume-hierarchy machinery into dosimetry.

\subsection{Projected area and the Cauchy bound}\label{sec:geom-cauchy}

Integrating $\Sab$ over the lit surface $\Sigma_+$ (where
$\mu\equiv\nhat\cdot(-\khat)>0$) gives total absorbed power as a single geometric
integral,
\begin{equation}\label{eq:Pabs-Aperp}
  P_{\mathrm{abs}}(\khat) = \Sinc\,T_0\int_{\Sigma_+}\!\mu\,\diff A
  = \Sinc\,T_0\,\Aperp(\khat),
  \qquad
  \Aperp(\khat) \equiv \int_{\Sigma_+}\mu\,\diff A .
\end{equation}
$\Aperp(\khat)$ is the projected (shadow) area of the body normal to $\khat$. For
a convex body it equals the area of the orthogonal projection. For a non-convex
body \cref{eq:Pabs-Aperp} overcounts unless self-occluded elements are removed,
that is unless $\mu$ is replaced by $\mu\,O$; with that correction the identity
holds for any body by a flux argument.

Averaging the projected area over all directions yields the classical Cauchy
result~\cite{Cauchy1841} and its non-convex generalisation through the exposure
fraction.

\begin{theorem}[Generalised Cauchy formula]\label{thm:gen-cauchy}
For any body with surface $\Sigma$ and exposure fraction $\eta$, the
direction-averaged absorbed power is
\begin{equation}\label{eq:gen-cauchy}
  \langle P_{\mathrm{abs}}\rangle
  = \frac{1}{4\pi}\int_{S^2} P_{\mathrm{abs}}(\khat)\,\diff\Omega
  = \Sinc\,T_0\,\frac{\Aab}{4},
  \qquad
  \Aab \equiv \int_\Sigma \eta(\rr)\,\diff A .
\end{equation}
\end{theorem}

\begin{proof}
Exchange the order of integration (Fubini, the integrand is non-negative) and use
the rotational symmetry of the inner integral,
$\int_{S^2}\ReLU[\nhat\cdot(-\khat)]\,\diff\Omega = \pi$, independent of $\nhat$.
With the occlusion factor restricting the surface integral to visible elements,
$\langle P_{\mathrm{abs}}\rangle = (\pi/4\pi)\,\Sinc\,T_0\,\Aab = \Sinc\,T_0\,\Aab/4$.
\end{proof}

The absorption area $\Aab\le A$ collapses all of self-shadowing into one scalar;
for convex bodies $\eta\equiv 1$ and $\Aab=A$, recovering $\langle\Aperp\rangle =
A/4$. Replacing $T_0$ with the exact averaged transmission $\Tavg$
(\cref{sec:fresnel}) makes \cref{eq:gen-cauchy} exact at any frequency.

\subsection{Principal curvatures and the Fock radius}\label{sec:geom-curvature}

The geometric law treats each surface element as locally flat. The wave
corrections do not: the curvature near the shadow edge sets the width of the
diffractive penumbra. The relevant quantity is not the mean curvature but the
normal curvature in the plane of incidence.

At a smooth surface point the shape operator has two real eigenvalues, the
principal curvatures $\kappa_1\ge\kappa_2$, with orthonormal principal directions
$\bm{e}_1,\bm{e}_2$ in the tangent plane (sign convention: $\kappa>0$ for an
outward-bulging convex surface). Their sum is twice the mean curvature,
$H=\kappa_1+\kappa_2$, the quantity that enters the physical-optics curvature
term $\Sab \propto \ReLU(\mu)\,[1+\mu H/k]$ used at the curvature level of the
grid (\cref{sec:grid}). AEGIS estimates $\kappa_1,\kappa_2,\bm{e}_1$ per triangle
from a local quadric (osculating-jet) fit over a centroid neighbourhood,
validated on analytic sphere and cylinder meshes to a few percent.

The normal curvature in an arbitrary tangent direction follows from Euler's
theorem. Let $\varphi$ be the angle between $\bm{e}_1$ and the incidence ray
projected into the tangent plane. Then
\begin{equation}\label{eq:euler}
  \kappa_t(\varphi) = \kappa_1\cos^2\varphi + \kappa_2\sin^2\varphi,
\end{equation}
and the in-incidence-plane radius of curvature is $\Rfock = 1/\kappa_t$. This
$\Rfock$, not $2/H$, is the radius the Fock detour parameter of
\cref{sec:fock-local} requires.

\begin{proposition}[The Fock radius is the in-plane normal curvature]\label{prop:fock-radius}
Let a ray strike a circular cylinder of radius $R$ normal to its axis. The
in-incidence-plane Fock radius is $\Rfock = R$, whereas the mean-curvature
surrogate gives $2/H = 2R$, twice too large. The two disagree by a factor of $2$,
and the diffractive penumbra width $\propto(kR)^{-1/3}$ they predict differs by
$2^{1/3}\approx 1.26$.
\end{proposition}

\begin{proof}
For a cylinder of radius $R$ the principal curvatures are $\kappa_1=1/R$ (around
the cross-section) and $\kappa_2=0$ (along the axis), so $H=\kappa_1+\kappa_2=1/R$
and $2/H=2R$. A ray normal to the axis projects, in the tangent plane, onto the
cross-sectional principal direction $\bm{e}_1$, so $\varphi=0$ and
\cref{eq:euler} gives $\kappa_t=\kappa_1=1/R$, hence $\Rfock=R$. The penumbra
width scales as $(k\Rfock)^{-1/3}$ (\cref{sec:fock-local}); substituting $2R$ for
$R$ rescales it by $2^{1/3}$.
\end{proof}

The cylinder is a primary diffraction oracle (\cref{sec:validation}), so the
factor-of-two error is not academic: it would corrupt the very case the gate is
validated against. For a sphere $\kappa_1=\kappa_2=1/R$ and \cref{eq:euler} gives
$\Rfock=R$ for every direction, as it must by isotropy. On a flat face
$\kappa_t\to 0$ and $\Rfock\to\infty$, which returns the gate to a sharp ReLU.
